\chapter{希尔伯特空间}\label{ch:b3:hilbert}

希尔伯特空间是范数来自某个
内积的巴拿赫空间 --- 这一额外结构使无限维中
几乎全部欧几里得几何得以恢复：
正交投影存在，每个连续泛函都是
与某固定向量的内积（里斯），标准正交基
能把每个向量展开为收敛级数，
并满足勾股型的平方和关系（帕塞瓦尔）。本章的高潮是兑现一笔
旧“欠账”：三角函数系是
$L^2$ 的标准正交基，故帕塞瓦尔恒等式对 \emph{每一个}
平方可积函数都成立 --- 这一结论本科二年级仅能
对分段 $\mathcal C^1$ 函数证明。本章最后
以拉克斯–米尔格拉姆引理收尾，它是变分方法
处理微分方程的主力引理。

全文设 $H$ 为 $K = \R$ 或 $\C$ 上的向量空间。

\section{内积；投影定理}

\begin{definition}\label{def:b3:hilbert:inner}
称 \emph{内积}\index{内积} 为映射 $\langle
\cdot,\cdot\rangle \colon H\times H \to K$，关于第二个
变元线性，且满足 $\langle y, x\rangle =
\overline{\langle x, y\rangle}$，且 $\langle x, x\rangle > 0$
当 $x \neq 0$ 时成立。它诱导出范数 $\norm x = \langle x,
x\rangle^{1/2}$，以及 \emph{柯西–施瓦茨不等式}
$\abs{\langle x, y\rangle} \leq \norm x\norm y$（本科二年级的
证明 --- 判别式法 --- 仍然成立），以及
\emph{平行四边形恒等式}
\[
\norm{x + y}^2 + \norm{x - y}^2 = 2\norm x^2 + 2\norm y^2 .
\]
称 \emph{希尔伯特空间}\index{希尔伯特空间} 为按该范数完备的
内积空间。例如 $\ell^2$
（\cref{pb:b3:banach:1}）及最基本的 $L^2(\mu)$
其内积为 $\langle f, g\rangle = \int\bar fg\,\dd\mu$ ---
由里斯–费舍尔定理（\cref{thm:b3:lp:complete}）知其完备；
内积的有限性由柯西–施瓦茨不等式推出（即 $= $ 赫尔德不等式，取
$p = q = 2$ 时）。
\end{definition}

\begin{theorem}[向闭凸集的投影]
\label{thm:b3:hilbert:projection}
设 $C \neq \varnothing$ 为闭 \emph{凸}子集，
它位于希尔伯特空间 $H$ 中，又设 $x \in H$。存在唯一的 $p_C(x)
\in C$，使得\index{投影定理}
\[
\norm{x - p_C(x)} = d(x, C),
\]
它由下列性质刻画：$\operatorname{Re}\langle x - p_C(x),\ c -
p_C(x)\rangle \leq 0$ 对所有 $c \in C$ 成立。映射 $p_C$ 是
$1$-利普希茨的。
\end{theorem}

\begin{proof}
令 $d = d(x, C)$，取 $(c_n) \subseteq C$ 使 $\norm{x - c_n}
\to d$。对 $x - c_n$ 与 $x - c_m$ 应用平行四边形恒等式：
\[
\norm{c_n - c_m}^2 = 2\norm{x - c_n}^2 + 2\norm{x - c_m}^2 -
4\,\bigl\|x - \tfrac{c_n + c_m}2\bigr\|^2
\leq 2\norm{x{-}c_n}^2 + 2\norm{x{-}c_m}^2 - 4d^2
\]
（由凸性，中点属于 $C$）：右端趋于
$0$，故 $(c_n)$ 为柯西列，其极限 $p \in C$（由闭性）
并取得 $d$。唯一性：两个极小点由同一
恒等式给出 $\norm{p - p'}^2 \leq 2d^2 + 2d^2 - 4d^2 = 0$。

刻画：对 $c \in C$，$t \in \intoc01$，向量
$p + t(c - p) \in C$，于是
\[
d^2 \leq \norm{x - p - t(c-p)}^2
= d^2 - 2t\operatorname{Re}\langle x - p, c - p\rangle +
t^2\norm{c-p}^2 ;
\]
两边除以公因子，再令 $t \to 0^+$：$\operatorname{Re}\langle x - p, c -
p\rangle \leq 0$。反之，由此不等式得 $\norm{x -
c}^2 = \norm{x - p}^2 - 2\operatorname{Re}\langle x - p, c -
p\rangle + \norm{p - c}^2 \geq \norm{x-p}^2$。利普希茨性：对
对 $x, y$，投影分别为 $p, q$，将两个变分
不等式（分别取 $c = q$ 与 $c = p$）相加：
$\operatorname{Re}\langle x - y - (p - q), p - q\rangle \geq
0$，于是 $\norm{p - q}^2 \leq \operatorname{Re}\langle x - y, p -
q\rangle \leq \norm{x - y}\norm{p - q}$。
\end{proof}

\begin{theorem}[正交分解]
\label{thm:b3:hilbert:decomposition}
设 $F$ 为 \emph{闭子空间}（位于 $H$ 中）。则 $p_F$ 是
线性的，$x - p_F(x) \perp F$ 对所有 $x$ 成立，且\index{正交
补}
\[
H = F \oplus F^\perp,
\qquad F^\perp = \{y : \langle y, f\rangle = 0\ \forall f\in
F\},
\qquad (F^\perp)^\perp = F .
\]
对一般子空间，$(F^\perp)^\perp = \bar F$；特别地，
$F$ 稠密当且仅当 $F^\perp = \{0\}$。
\end{theorem}

\begin{proof}
对子空间，利用变分刻画，取 $c =
p_F(x) \pm f$（$f \in F$，两种符号，以及复数情形的 $\iu f$
）推出 $\langle x - p_F(x), f\rangle = 0$：残差
与 $F$ 正交。分解 $x = p_F(x) + (x
- p_F(x))$，且 $F \cap F^\perp = \{0\}$（$\langle y, y
\rangle = 0$）；$p_F$ 的线性由此类分解的
唯一性推出（两侧关于它们均线性）。$(F^\perp)
^\perp \supseteq F$ 总成立；反之若 $x \perp F^\perp$，
记 $x = f + g$：$g = x - f \in F^\perp$，且 $\langle g,
g\rangle = \langle x, g\rangle - \langle f, g\rangle = 0$：$x
= f \in F$。对一般子空间 $F$：$F^\perp = \bar
F^{\,\perp}$（内积的连续性），于是
$(F^\perp)^\perp = \bar F$（由闭情形）；稠密性
即 $\bar F = H$，当且仅当 $F^\perp = 0$。
\end{proof}

\begin{theorem}[里斯表示定理]\label{thm:b3:hilbert:riesz}
对每个连续线性泛函 $\varphi \in H'$，
存在唯一的 $a \in H$，使\index{里斯表示定理}
\[
\varphi(x) = \langle a, x\rangle \quad (x \in H),
\qquad \norm\varphi_{H'} = \norm a .
\]
\end{theorem}

\begin{proof}
若 $\varphi = 0$，则 $a = 0$。否则 $F = \ker\varphi$ 为
闭的真子空间；取 $u \in F^\perp$，$\norm u = 1$
（由\cref{thm:b3:hilbert:decomposition}：$F^\perp \neq 0$，
因 $F \neq H$）。对任意 $x$，向量 $\varphi(x)u -
\varphi(u)x \in \ker\varphi$，从而 $\perp u$：
\[
0 = \langle u, \varphi(x)u - \varphi(u)x\rangle
= \varphi(x) - \varphi(u)\langle u, x\rangle :
\qquad \varphi(x) = \langle
\overline{\varphi(u)}\,u,\ x\rangle .
\]
于是 $a = \overline{\varphi(u)}u$ 满足要求。唯一性：$\langle a
- a', x\rangle = 0$ 对所有 $x$ 成立，取 $x = a - a'$。范数：
$\abs{\varphi(x)} \leq \norm a\norm x$（柯西–施瓦茨不等式），等号在
$x = a$ 处成立。
\end{proof}

\begin{example}[一个投影的完整计算]
\label{ex:b3:hilbert:projectionexample}
在 $H = L^2(\intcc01)$ 中，用仿射函数对
$f(x) = x^2$ 作最佳逼近是什么？子空间 $F =
\operatorname{Vect}(1, x)$ 是闭的（有限维），
且 $p_F(f) = a + bx$ 由残差分别与
$1$ 和 $x$ 正交所刻画：
\[
\int_0^1(x^2 - a - bx)\,\dd x = 0,
\qquad
\int_0^1x\,(x^2 - a - bx)\,\dd x = 0,
\]
即 $\frac13 = a + \frac b2$，且 $\frac14 = \frac a2 +
\frac b3$：$a = -\frac16$，$b = 1$。于是 $p_F(x^2) = x -
\frac16$，误差为
\[
d(f, F)^2 = \int_0^1\Bigl(x^2 - x + \frac16\Bigr)^2\dd x =
\frac1{180},
\qquad d(f, F) = \frac1{6\sqrt5} .
\]
两点值得留意。第一，这一计算不过是
一个 $2\times2$ 线性方程组 --- \emph{正规
方程}；对单项式基，其矩阵
$\bigl(\frac1{i+j+1}\bigr)$ 是著名的
病态希尔伯特矩阵，而先行正交化
（勒让德多项式，\cref{pb:b3:hilbert:1}）即是对策。
第二，用仿射函数对 $x^2$ 作最佳\emph{一致}逼近
则不同，其结果为（$x - \frac18$，
由等振荡得出）：每种范数各有其几何，唯有
希尔伯特型范数才归结为一个线性方程组。
\end{example}

\section{标准正交基}

\begin{definition}\label{def:b3:hilbert:onb}
称族 $(e_i)_{i\in I}$ 为\emph{标准正交的}，若 $\langle
e_i, e_j\rangle = \delta_{ij}$，并称\emph{希尔伯特
基}\index{希尔伯特基}（标准正交基）若另外
其有限线性组合在 $H$ 中稠密（即该族是
\emph{完全}的）。下面只讨论可数情形 $I = \N$，由
格拉姆–施密特方法即覆盖每个\emph{可分}的 $H$
（\cref{prop:b3:hilbert:gramschmidt}）。
\end{definition}

\begin{theorem}[贝塞尔、帕塞瓦尔]\label{thm:b3:hilbert:parseval}
设 $(e_n)_{n\in\N}$ 为 $H$ 中的标准正交族，且 $c_n(x) =
\langle e_n, x\rangle$。
\begin{enumerate}
  \item （贝塞尔）$\sum_n\abs{c_n(x)}^2 \leq \norm x^2$，且
        级数 $\sum_nc_n(x)e_n$ 在 $H$ 中收敛，其和为
        $p_F(x)$，$F = \overline{\operatorname{Vect}}(e_n)$。
  \item 下列各条件等价：（i）$(e_n)$ 为希尔伯特
        基；（ii）$x = \sum_nc_n(x)e_n$ 对每个 $x$ 成立；（iii）
        \emph{帕塞瓦尔恒等式}\index{帕塞瓦尔恒等式}：$\norm x^2
        = \sum_n\abs{c_n(x)}^2$ 对每个 $x$ 成立；（iv）与所有 $e_n$ 正交的
        向量只有 $0$。
  \item 若 $(e_n)$ 为希尔伯特基，则 $x \mapsto (c_n(x))_n$
        是 $H \to \ell^2$ 的等距同构
        （\emph{每个}无穷维可分希尔伯特空间
        都“是” $\ell^2$），且 $\langle x, y\rangle =
        \sum_n\overline{c_n(x)}c_n(y)$。
\end{enumerate}
\end{theorem}

\begin{proof}
（1）对有限的 $N$，$x - \sum_{n\leq N}c_ne_n \perp e_k$（$k
\leq N$），由毕达哥拉斯定理得 $\norm x^2 = \sum_{n\leq
N}\abs{c_n}^2 + \norm{x - \sum_{n\leq N}c_ne_n}^2$，即贝塞尔不等式。
部分和 $S_N = \sum_{n\leq N}c_ne_n$ 为柯西列：
$\norm{S_N - S_M}^2 = \sum_{M<n\leq N}\abs{c_n}^2$，它是一个
收敛级数的尾部；极限属于 $F$，且 $x - \lim S_N
\perp$ 与每个 $e_k$ 正交（由连续性），从而 $\perp F$。由正交分解的唯一性，
$\lim S_N = p_F(x)$。

（2）（i）$\Rightarrow$（ii）：$F = H$，故 $p_F = \mathrm{id}$。
（ii）$\Rightarrow$（iii）：取极限后用毕达哥拉斯定理（$\norm{S_N}^2
= \sum_{n \leq N}\abs{c_n}^2 \to \norm x^2$）。
（iii）$\Rightarrow$（iv）：$x \perp$ 所有 $e_n$ 给出 $\norm x^2 =
0$。（iv）$\Rightarrow$（i）：$F^\perp = \{0\}$
（与所有 $e_n$ 正交即与 $F$ 正交），故 $F$
稠密（由\cref{thm:b3:hilbert:decomposition}）；但 $F$，作为
闭包本身已是闭的：$F = H$。

（3）该映射是线性的，由（iii）知其保距（从而为单射），
并且满射：给定 $(c_n) \in \ell^2$，级数 $\sum
c_ne_n$ 收敛（同（1）一样，由柯西准则）到一个原像。内积
公式由（iii）的极化恒等式得出，或由直接的极限
计算得到。
\end{proof}

\begin{proposition}[格拉姆--施密特]\label{prop:b3:hilbert:gramschmidt}
设 $(x_n)$ 为线性无关序列。
归纳地令 $\tilde e_n = x_n - \sum_{k<n}\langle e_k,
x_n\rangle e_k$ 且 $e_n = \tilde e_n/\norm{\tilde e_n}$
得到标准正交的 $(e_n)$，且有相同的有限张成：
$\operatorname{Vect}(e_1, \dots, e_n) = \operatorname{Vect}
(x_1, \dots, x_n)$。因此每个可分希尔伯特空间
（即含可数稠密子集者）都有希尔伯特
基。\index{格拉姆--施密特正交化}
\end{proposition}

\begin{proof}
归纳法：由构造知 $\tilde e_n \perp e_k$（$k < n$），
又由线性无关性得 $\tilde e_n \neq 0$；张成在每一步
均相同（三角形式的基变换）。对可分的 $H$：
从一个稠密序列中取出线性无关的
子族，使其张成稠密（舍去每个落在其前面向量张成中的向量，
张成不变），再标准正交化：
所得族是完全的。
\end{proof}

\begin{theorem}[三角函数系；终于证明帕塞瓦尔恒等式]
\label{thm:b3:hilbert:fourier}
在 $L^2(\intcc{-\pi}\pi)$ 中，取 $\langle f, g\rangle =
\frac1{2\pi}\int_{-\pi}^\pi \bar fg$，则族 $e_n(t) =
\eu^{\iu nt}$，$n \in \Z$，是希尔伯特基。因此，
对\emph{每一个} $f \in L^2$ --- 特别地，每个分段
连续的 $2\pi$ 周期函数 $f$ ---，其中 $c_n(f) =
\frac1{2\pi}\int_{-\pi}^{\pi}f(t)\eu^{-\iu nt}\dd t$:
\[
f = \sum_{n\in\Z}c_n(f)\,\eu^{\iu nt} \ \ \text{在 } L^2,
\qquad
\frac1{2\pi}\int_{-\pi}^{\pi}\abs f^2 =
\sum_{n\in\Z}\abs{c_n(f)}^2 .
\]
这在完全一般的意义下证明了
本科二年级在本阶段不加证明地承认的帕塞瓦尔恒等式。\index{帕塞瓦尔恒等式（一般形式）}
\end{theorem}

\begin{proof}
标准正交性是直接计算（本科二年级）。完全性：设
$f \in L^2$ 且 $\perp$ 所有 $e_n$，即所有傅里叶
系数全为零。连续的 $2\pi$ 周期函数
在 $L^2(\intcc{-\pi}\pi)$ 中稠密：事实上 $\mathcal
C_c(\intoo{-\pi}\pi)$ 是稠密的
（见\cref{thm:b3:lp:density}(2)），且这类函数可以
周期且连续地延拓。三角多项式
在连续周期函数中按 $\norm\cdot_\infty$ 稠密
（斯通--魏尔斯特拉斯，见\cref{cor:b3:complete:weierstrass}(c)），
且 $\norm\cdot_2 \leq \norm\cdot_\infty$，故三角
多项式在 $L^2$ 中稠密。但 $f \perp$ 每一个
三角多项式，因而 $f \perp$ 某一稠密子空间：$f
\in (\text{稠密})^\perp = \{0\}$
（见\cref{thm:b3:hilbert:decomposition}）。由
\cref{thm:b3:hilbert:parseval}的准则（iv）即得结论；（ii）与（iii）
展开为上式（对可数集 $\Z$ 重新编号；
双端级数无条件收敛 --- 任一穷竭族上的部分和
都收敛，由 $\ell^2$ 尾项
论证）。
\end{proof}

\begin{theorem}[拉克斯--米尔格拉姆]\label{thm:b3:hilbert:laxmilgram}
设 $H$ 为实希尔伯特空间，$a \colon H\times H \to
\R$ 为双线性、\emph{连续}的（$\abs{a(u,v)} \leq M\norm
u\norm v$），且\emph{强制的}（$a(u, u) \geq \alpha\norm u^2$，
$\alpha > 0$）。则对任意 $\varphi \in H'$，存在一个
唯一的 $u \in H$ 使得\index{拉克斯--米尔格拉姆定理}
\[
a(u, v) = \varphi(v) \qquad \text{对所有 } v \in H .
\]
\end{theorem}

\begin{proof}
固定 $u$，$v \mapsto a(u, v)$ 是连续线性的：由里斯
表示定理，存在唯一的 $Au \in H$ 使 $a(u,v) = \langle Au,
v\rangle$；$A$ 线性，且 $\norm{Au} \leq M\norm u$
（表示元唯一，再估界）。强制性：
$\alpha\norm u^2 \leq a(u,u) = \langle Au, u\rangle \leq
\norm{Au}\norm u$，故 $\norm{Au} \geq \alpha\norm u$：$A$ 是
单射，且值域闭（像列 $Au_n$ 为柯西列
则 $u_n$ 亦为柯西列）。值域稠密：若 $w \perp
\operatorname{im}A$，则 $0 = \langle Aw, w\rangle \geq
\alpha\norm w^2$。既闭又稠，故 $A$ 是双射。给定
$\varphi$，取 $f$ 表示它（由里斯表示定理），并令 $u = A^{-1}f$：
$a(u, v) = \langle f, v\rangle = \varphi(v)$，且
唯一（由 $a(u - u', \cdot) = 0$ 及强制性）。
\end{proof}

\begin{remark}\label{rem:b3:hilbert:variational}
当 $a$ 对称时，拉克斯--米尔格拉姆解是
\emph{能量}泛函 $J(v) = \frac12a(v,v) -
\varphi(v)$ 的唯一极小元（见\cref{exo:b3:hilbert:9}）：
一举得到变分问题解的存在性。将其应用于合适的
函数空间（后续课程中的索伯列夫空间），这
便可求解微分方程的边值问题 ---
正是通往偏微分方程的现代入口。
\end{remark}

\section{练习}

\begin{exercise}[$\star$]\label{exo:b3:hilbert:1}
(a) 证明极化恒等式（实情形：$4\langle x,
y\rangle = \norm{x+y}^2 - \norm{x-y}^2$；复情形：
四项版本）。
(b) 证明 $\norm\cdot_1$（在 $L^1(\intcc01)$ 上）及
$\norm\cdot_\infty$（在 $\mathcal C(\intcc01)$ 上）都违反
平行四边形法则：这些范数都不来自任何内积。
\end{exercise}

\begin{exercise}[$\star$]\label{exo:b3:hilbert:2}
在实希尔伯特空间 $H = L^2(\intcc01)$ 中：
(a) 求 $f$ 在常值函数构成的子空间上的
投影，并解释其意义；
(b) 求到 $\{g : g = 0 \text{ 几乎处处于 }
\intcc0{1/2}\}$ 上的投影；
(c) 计算 $d\bigl(x \mapsto x,\ \operatorname{Vect}(\mathbf
1)\bigr)$。
\end{exercise}

\begin{exercise}[$\star\star$]\label{exo:b3:hilbert:3}
(a) 证明：对子空间 $F$，$F$ 稠密 $\iff$ $F^\perp =
\{0\}$，并在 $\ell^2$ 中给出一个\emph{真}
稠密子空间的例子（这表明 $F^\perp = 0$ 并不蕴含 $F = H$：
分解定理确实需要 $F$ 闭）。
(b) 证明：若在范数意义下 $x_n \to x$ 且 $y_n \to y$，则
$\langle x_n, y_n\rangle \to \langle x, y\rangle$，并指出
本章中两处用到这一连续性的地方。
\end{exercise}

\begin{exercise}[$\star\star$]\label{exo:b3:hilbert:4}
对 $1, x, x^2$ 在 $L^2(\intcc{-1}1)$ 中作格拉姆–施密特正交化
（勒贝格测度）：求出前三个标准化的
\emph{勒让德多项式}，并验证它们与
$\sqrt{n + \frac12}\,P_n$ 相符，其中 $P_n$ 为
\cref{pb:b3:hilbert:1} 中的罗德里格斯多项式。
\end{exercise}

\begin{exercise}[$\star\star$]\label{exo:b3:hilbert:5}
对 $f(t) = t$ 应用帕塞瓦尔恒等式（\cref{thm:b3:hilbert:fourier}），
及 $f(t) = t^2$ 在 $\intcc{-\pi}\pi$ 上 --- 现在对这两个函数可以合法地应用
（它们是连续的，但此前恒等式需要
对分段 $\mathcal C^1$ 函数在绕回处的
间断点作细致处理），由此推出
\[
\sum_{n\geq1}\frac1{n^2} = \frac{\pi^2}6,
\qquad
\sum_{n\geq1}\frac1{n^4} = \frac{\pi^4}{90} .
\]
\end{exercise}

\begin{exercise}[$\star\star$]\label{exo:b3:hilbert:6}
(a) 求 $a \in L^2(\intcc01)$，使得 $\int_0^{1/2}f =
\langle a, f\rangle$ 对所有 $f$ 成立，并计算
该泛函的 $\norm\varphi$。
(b) 证明：求值映射 $f \mapsto f(\frac12)$，它定义在子空间
$\mathcal C(\intcc01) \subseteq
L^2(\intcc01)$ 上，它\emph{不}连续（关于
$\norm\cdot_2$ 而言），于是不存在里斯表示元（求值
并不是 $L^2$ 意义下的概念）。
\end{exercise}

\begin{exercise}[$\star\star\star$]\label{exo:b3:hilbert:7}
设 $H$ 可分，$(e_n)$ 为其希尔伯特基，$(x_k)$ 为
一个有界序列。
(a) 证明：某个子列\emph{弱收敛}：存在
$x$，使得 $\langle y, x_{k_j}\rangle \to \langle y,
x\rangle$ 对每个 $y \in H$ 成立。\emph{（对
系数 $\langle e_n, x_k\rangle$ 作对角线法抽取；借助
贝塞尔不等式与范数的一致有界性组装出 $x$。）}
(b) 证明 $e_n \rightharpoonup 0$ 但 $\norm{e_n} = 1$：弱
极限可能丢失范数。证明 $\norm x \leq
\liminf\norm{x_{k_j}}$ 在 (a) 中成立。
\end{exercise}

\begin{exercise}[$\star\star$]\label{exo:b3:hilbert:8}
（伴随算子）对 $T \in \mathcal L(H)$，证明存在唯一的
$T^* \in \mathcal L(H)$，使得 $\langle Tx, y\rangle = \langle
x, T^*y\rangle$（由里斯表示定理），并且 $\vertiii{T^*} = \vertiii T$。
求移位算子 $S$ 在 $\ell^2$ 上的伴随，并证明
$\ker T^* = (\operatorname{im}T)^\perp$ --- 由此推出
$\overline{\operatorname{im}T} = (\ker T^*)^\perp$。
\end{exercise}

\begin{exercise}[$\star\star$]\label{exo:b3:hilbert:9}
设 $a$ 如拉克斯–米尔格拉姆定理所设，且 \emph{对称}。
证明：$u$ 满足 $a(u, \cdot) = \varphi$ 当且仅当 $u$ 使
$J(v) = \frac12a(v, v) - \varphi(v)$ 取得最小值，并且最小值
恰在一点处取得。\emph{（配方：
$J(u + w) - J(u) = \frac12a(w,w) \geq
\frac\alpha2\norm w^2$。）} 应用：由拉克斯–米尔格拉姆定理
重新导出闭子空间上的投影定理。
\end{exercise}

\begin{exercise}[$\star\star\star$]\label{exo:b3:hilbert:10}
（哈尔系）在 $\intcc01$ 上，令 $h_{0} = \mathbf 1$，并且
对 $n = 2^j + k$（$j \geq 0$，$0 \leq k < 2^j$）：
\[
h_n = 2^{j/2}\Bigl(\mathbf 1_{[k2^{-j},\,(k +
\frac12)2^{-j})} - \mathbf 1_{[(k+\frac12)2^{-j},\,(k+1)2^{-j})}
\Bigr).
\]
证明 $(h_n)_{n\geq0}$ 在 $L^2(\intcc01)$ 中是标准正交的，
并且完全。\emph{（正交性：支集不交或彼此嵌套；
完全性：有限张成包含所有二进阶梯函数，
而它们稠密（由 \cref{thm:b3:lp:density}(1) 及区间的二进
逼近得到。）} 哈尔系是
小波的先驱。
\end{exercise}

\begin{exercise}[$\star\star$]\label{exo:b3:hilbert:11}
（正交投影的刻画）设 $H$ 为希尔伯特
空间，且 $P \in \mathcal L(H)$ 满足 $P^2 = P$，$P \neq 0$。
证明下列三条件等价：
(i) $P$ 为到
$\operatorname{im}P$ 上的正交投影；
(ii) $P = P^*$（见\cref{exo:b3:hilbert:8}）；
(iii) $\vertiii P = 1$。
\emph{（对 (iii) $\Rightarrow$ (i)：若某 $x \in
(\ker P)^\perp$ 满足 $Px \neq x$，考虑 $x + t(Px - x)$
--- 或者直接：对 $u \in \operatorname{im}P$ 及 $v \in
\ker P$，展开 $\norm{P(u + tv)}^2 \leq \norm{u + tv}^2$
对所有 $t \in \R$ 成立，并推出 $\langle u, v\rangle = 0$。）}
给出 $\R^2$ 上一个非正交投影，并计算
其范数。
\end{exercise}

\begin{exercise}[$\star\star\star$]\label{exo:b3:hilbert:12}
（冯·诺伊曼遍历定理）设 $U \in \mathcal L(H)$ 为
\emph{酉算子}（$U^*U = UU^* = I$），记 $F = \ker(U - I)$
为不动点空间，$P$ 为到 $F$ 上的正交投影，并记
$A_n = \frac1n\sum_{k=0}^{n-1}U^k$.
(a) 证明 $\ker(U - I) = \ker(U^* - I)$ \emph{（由
$\norm{Ux - x}^2 = 2\norm x^2 - 2\operatorname{Re}\langle
Ux, x\rangle$ 及酉性）}，并推出
$\overline{\operatorname{im}(U - I)} = F^\perp$.
(b) 证明 $A_nx \to x$ 对 $x \in F$ 成立，且 $A_nx \to 0$
对 $x \in \operatorname{im}(U - I)$ 成立 \emph{（裂项相消）}，
进而对 $x \in \overline{\operatorname{im}(U - I)}$ 成立，
（利用一致有界性 $\vertiii{A_n} \leq 1$）。
(c) 综上，$A_nx \to Px$ 对\emph{每个} $x \in H$ 成立 ---
即时间平均收敛于到不变向量上的投影。
(d) 具体地，对 $H = L^2(\R/\Z)$ 及 $Uf = f(\cdot +
\alpha)$ 其中 $\alpha$ 为无理数，确定 $F$（利用
傅里叶级数，见\cref{thm:b3:hilbert:fourier}），并推出
$\frac1n\sum_{k<n}f(x + k\alpha) \to \int_0^1f$ 在
$L^2$ 意义下成立：即无理旋转的 $L^2$ 等分布性。
\end{exercise}

\section{问题：正交多项式}

\begin{problem}[{周末问题 --- 勒让德、埃尔米特与
高斯求积}]\label{pb:b3:hilbert:1}
设 $I \subseteq \R$ 为区间，$w > 0$ 为连续
\emph{权函数}在 $I$ 的内部，且使 $\int_I
\abs t^nw(t)\dd t < \infty$ 对所有 $n$ 成立；考虑 $H = L^2(I,
w\,\dd\lambda)$，其内积为 $\langle f, g\rangle = \int_I \bar
fg\,w$。对 $1, t, t^2, \dots$ 施行格拉姆–施密特正交化，
得到 \emph{正交多项式} $(p_n)$（关于 $w$，
首一规范化：$p_n = t^n + \cdots$）。\index{正交
多项式}

\medskip
\noindent\textbf{第一部分 --- 一般理论。}
\begin{enumerate}
  \item 证明 $p_n$ 与每个次数
        $< n$ 的多项式正交，且 $(p_0, \dots, p_n)$ 是
        $\R_n[t]$ 的一组基。
  \item （三项递推关系）证明存在实数 $a_n,
        b_n$ 使得
        \[
        p_{n+1}(t) = (t - a_n)\,p_n(t) - b_n\,p_{n-1}(t),
        \qquad b_n = \frac{\norm{p_n}^2}{\norm{p_{n-1}}^2} >
        0 .
        \]
        \emph{（将 $t\,p_n$ 在基 $(p_k)_{k \leq
        n+1}$ 下展开，并由正交性消去系数，其中用到
        $\langle tp_n, p_k\rangle = \langle p_n,
        tp_k\rangle$.)}
  \item （根）证明 $p_n$ 有 $n$ 个\emph{互异}的根，
        全部位于 $I$ 的内部。\emph{（设 $t_1 < \dots < t_m$ 为
        $p_n$ 的内部变号点，若 $m < n$，则
        考察 $p_n$ 与 $\prod_{i\leq m}(t - t_i)$ 的内积，
        与正交性矛盾。）}
\end{enumerate}

\medskip
\noindent\textbf{第二部分 --- 勒让德（$I = \intcc{-1}1$，$w =
1$）。}定义 $P_n(t) = \frac{1}{2^nn!}\,\frac{\dd^n}{\dd
t^n}\bigl[(t^2 - 1)^n\bigr]$（罗德里格斯公式）。
\begin{enumerate}[resume]
  \item 证明 $\deg P_n = n$，且首项系数为
        $\frac{(2n)!}{2^n(n!)^2}$，并分部积分
        $n$ 次，可得 $\langle P_n, Q\rangle = 0$，对任一
        多项式 $Q$（次数 $< n$）成立；于是 $P_n$（相差
        规范化常数）即为关于 $w = 1$ 的正交多项式。
  \item 计算 $\norm{P_n}_2^2 = \frac{2}{2n+1}$
        \emph{（对自身分部积分 $n$ 次，
        并化为贝塔/沃利斯积分，
        见\cref{exo:b3:product:8}）}。
  \item 证明归一化勒让德多项式构成
        $L^2(\intcc{-1}1)$ 的希尔伯特基
        \emph{（魏尔斯特拉斯，见\cref{cor:b3:complete:weierstrass}，
        以及 $\mathcal C$ 在 $L^2$ 中的稠密性）}，并
        将 $f(t) = \abs t$ 展开到 $2$ 次，求
        $L^2$ 意义下的最佳二次逼近 $\abs t$。
\end{enumerate}

\medskip
\noindent\textbf{第三部分 --- 埃尔米特（$I = \R$，$w(t) =
\eu^{-t^2}$）。}定义 $H_n(t) =
(-1)^n\eu^{t^2}\frac{\dd^n}{\dd t^n}\eu^{-t^2}$.
\begin{enumerate}[resume]
  \item 证明 $H_n$ 是 $n$ 次多项式，且首项系数为
        $2^n$，并证明 $H_{n+1} = 2tH_n -
        H_n'$，以及 $\langle H_m, H_n\rangle_w =
        \delta_{mn}\,2^nn!\sqrt\pi$ \emph{（再次分部积分）}。
  \item 证明埃尔米特族在 $L^2(\R,
        \eu^{-t^2}\dd t)$ 中完全，并允许使用
        \cref{ch:b3:fouriertransform}中的一个结论：若 $g \in L^1(\R)$ 满足
        $\int g(t)\eu^{-\iu\xi t}\dd t = 0$ 对所有 $\xi$ 成立，
        则 $g = 0$ 几乎处处成立。\emph{（设 $f \perp$ 所有 $H_n$，
        即 $\perp$ 所有多项式：证明 $z \mapsto \int
        f(t)\eu^{-t^2}\eu^{-\iu zt}\dd t$ 有良定义，
        将指数展开为级数，并论证求和与积分的
        交换可由控制收敛给出，并推出
        $f\eu^{-t^2}$ 的傅里叶变换为零。）}
\end{enumerate}

\medskip
\noindent\textbf{第四部分 --- 高斯求积。}固定 $n$，设
$t_1 < \dots < t_n$ 为 $p_n$ 的根（第一部分），并定义
求积系数 $w_i = \int_I \ell_i(t)\,w(t)\dd t$，其中
$\ell_i$ 为关于
$t_i$ 的拉格朗日插值基多项式。
\begin{enumerate}[resume]
  \item 证明求积公式 $Q(f) = \sum_iw_if(t_i)$
        对所有次数 $\leq n - 1$ 的多项式精确成立
        （由插值），事实上 --- 这是奇迹 --- 对
        所有次数 $\leq 2n - 1$ 的多项式也成立：记 $P =
        qp_n + r$，并对商 $q$ 利用正交性。
  \item 证明求积系数为正
        \emph{（将该公式用于 $\ell_i^2$，其次数为 $2n -
        2$）}，并由波利亚定理
        （见\cref{exo:b3:banach:9}）推出高斯求积收敛，即
        $Q_n(f) \to \int_I fw$ 对每个连续
        $f$ 在紧集 $I$ 上成立。
  \item 对 $n = 2$、$I = \intcc{-1}1$、$w = 1$，求
        节点 $\pm\frac1{\sqrt3}$ 与求积系数 $1, 1$，
        手算验证其对 $1, t, t^2, t^3$ 精确，并与
        在相同的两个求值点上的梯形公式
        作比较。
\end{enumerate}

\medskip
\noindent\textbf{第 V 部分 --- 切比雪夫：
振荡最优的多项式。}现取 $I = \intcc{-1}1$，且 $w(t) =
\frac1{\sqrt{1 - t^2}}$。
\begin{enumerate}[resume]
  \item 证明 $T_n(\cos\theta) = \cos n\theta$ 定义了一个
        多项式 $T_n$（次数为 $n$，由三角恒等式建立 $T_{n+1} =
        2t\,T_n - T_{n-1}$），
        首项系数为 $2^{n-1}$（$n \geq 1$）；
        并且由替换 $t = \cos\theta$ 可得
        \[
        \langle T_m, T_n\rangle_w =
        \int_0^\pi\cos m\theta\,\cos n\theta\,\dd\theta
        = 0 \ (m \neq n), \qquad
        \norm{T_0}_w^2 = \pi,\ \ \norm{T_n}_w^2 = \frac\pi2 :
        \]
        $T_n$ 即对于此
        权的正交多项式，而切比雪夫展开\emph{实为}傅里叶
        余弦级数的伪装形式。
  \item 显式求出 $n$ 个根 $t_k =
        \cos\frac{(2k-1)\pi}{2n}$，以及 $n + 1$ 个极值点
        $s_j = \cos\frac{j\pi}n$，它们是 $T_n$ 在 $\intcc{-1}1$ 上的
        极值点，其中 $T_n(s_j) = (-1)^j$，图像
        在 $\pm1$ 之间\emph{等振荡}。
  \item （极小极大性）证明：在所有\emph{首一}
        $n$ 次多项式中，多项式
        $2^{1-n}T_n$ 的上确界范数在
        $\intcc{-1}1$ 上最小，为 $2^{1-n}$，且它是
        唯一的极小化元。\emph{（若某首一 $P$ 满足
        $\sup\abs P < 2^{1-n}$，则差 $2^{1-n}T_n -
        P$ 的次数 $\leq n-1$，将在
        $n+1$ 个等振荡点上符号交替。）}
  \item 插值的应用：对节点 $t_1 < \dots
        < t_n$ 位于 $\intcc{-1}1$ 中，则
        $\mathcal C^n$ 函数的拉格朗日插值误差涉及
        $\omega(t) = \prod_i(t - t_i)$。证明：取切比雪夫根为节点
        可极小化
        $\sup_{\intcc{-1}1}\abs\omega$，并给出由此得到的界
        $\norm{f -
        L_nf}_\infty \leq \frac{\norm{f^{(n)}}_\infty}
        {2^{n-1}\,n!}$ --- 与等距节点相比较
        （陈述龙格现象，以之为警示）。
  \item 验证 $\abs{T_n'(\pm1)} = n^2$ \emph{（对
        $T_n(\cos\theta) = \cos n\theta$ 求导，并取极限
        $\theta \to 0, \pi$）}：以 $1$ 为界的多项式
        在 $\intcc{-1}1$ 上，其端点处的导数可大至 $n^2$
        （马尔可夫不等式表明不会更大 ---
        仅作陈述）。区间中哪一部分的导数界
        仅为 $O(n)$？
  \item （切比雪夫--高斯求积）证明：
        对权 $w$ 在 $n$ 个切比雪夫根处的高斯求积公式
        具有\emph{相等}的系数 $w_i = \frac\pi n$
        \emph{（对 $T_0, \dots, T_{n-1}$ 精确，外加
        三角和
        $\sum_{k=1}^n\cos\bigl(j\tfrac{(2k-1)\pi}{2n}\bigr) =
        0$ 对 $1 \leq j \leq n - 1$ 成立）}：在所有求积公式中最
        均匀的。对 $n = 3$ 写出该公式。
\end{enumerate}

\medskip
\noindent\textbf{第 VI 部分 --- 克里斯托费尔--达布，
交错性与雅可比矩阵。}回到一般的权；
$h_k = \norm{p_k}^2$（$p_k$ 为首一多项式），$b_k = h_k/h_{k-1}$。
\begin{enumerate}[resume]
  \item （最小范数）证明：在所有\emph{首一}
        $n$ 次多项式中，正交多项式 $p_n$
        是唯一使 $L^2(w)$ 范数最小者，并将该极小问题
        视为到
        $\R_{n-1}[t]$ 上的正交投影（见\cref{thm:b3:hilbert:projection}，或
        本科二年级所学的有限维投影）。
        第 14 题的极小极大性质即同一陈述，
        只需以 $L^\infty$ 代替 $L^2$：同一主角，两种
        范数。
  \item （克里斯托费尔--达布）对 $n$ 用归纳法
        并利用三项递推关系，证明恒等式
        \[
        \sum_{k=0}^{n}\frac{p_k(x)\,p_k(y)}{h_k}
        = \frac{p_{n+1}(x)\,p_n(y) -
        p_n(x)\,p_{n+1}(y)}{h_n\,(x - y)}
        \qquad (x \neq y),
        \]
        及其汇合形式（$y \to x$）：
        $\sum_{k\leq n}\frac{p_k(x)^2}{h_k} =
        \frac{p_{n+1}'(x)p_n(x) - p_n'(x)p_{n+1}(x)}{h_n}$。
  \item 由此推出 $p_n$ 与 $p_{n+1}$ 没有公共根，
        且对每个根 $x_0$（属于 $p_{n+1}$）：
        $p_n(x_0)\,p_{n+1}'(x_0) > 0$。由此得出根的
        \emph{交错性}：在 $p_{n+1}$ 的两个相邻根之间
        恰有 $p_n$ 的一个根。
  \item （雅可比矩阵）设 $J_n$ 为 $n\times n$
        对称三对角矩阵，主对角线元素为 $a_0,
        \dots, a_{n-1}$，次对角元为 $\sqrt{b_1},
        \dots, \sqrt{b_{n-1}}$。用归纳法证明
        $\det(tI_n - J_n) = p_n(t)$，故 $p_n$ 的根
        即某实对称矩阵的特征值 ---
        由此一行即可重证其为实数，并且（结合上述
        交错性），将正交多项式与
        \cref{ch:b3:spectral}中的谱理论联系起来。
  \item （综合）汇总三大经典族
        （勒让德、埃尔米特、切比雪夫）的对照表：
        区间、权、定义式、三项
        递推式、范数，以及各自的自然应用场合
        （紧集上的求积与逼近；高斯
        分析；极小极大与傅里叶余弦方法）。
        用一句话说明一般理论（第 I、VI 部分）
        给出了任何单独计算都无法得到的结论。
\end{enumerate}

\medskip
\noindent\textbf{第 VII 部分 --- 误差项，以及权背后的
核。}此处 $I$ 为紧集，且 $f \in \mathcal
C^{2n}(I)$。
\begin{enumerate}[resume]
  \item （高斯误差公式）设 $Hf$ 为
        次数 $\leq 2n - 1$ 且与 $f$ 匹配的插值多项式，
        并与 $f'$ 在节点 $t_1, \dots, t_n$ 处匹配（证明其存在性，
        及逐点误差
        \[
        f(t) - Hf(t) =
        \frac{f^{(2n)}(\xi_t)}{(2n)!}\;p_n(t)^2
        \]
        由通常的辅助函数论证）。通过将该恒等式
        对 $w$ 积分，并利用
        $f^{(2n)}$ 的极值进行夹逼，可得
        \[
        \int_I f\,w - Q_n(f)
        = \frac{f^{(2n)}(\xi)}{(2n)!}\,h_n
        \qquad\text{ 对某个 } \xi \in I,
        \]
        其中 $h_n = \norm{p_n}^2$（如第 VI 部分所定义）：高斯
        求积的误差由一个 $2n$ 阶导数给出，并以
        首一正交多项式范数的平方加权。
  \item （权即为克里斯托费尔函数值）利用
        再生核 $K_n(x, y) = \sum_{k=0}^{n-1}
        \frac{p_k(x)p_k(y)}{h_k}$（属于 $\R_{n-1}[t]$），以及
        $Q_n$ 直至次数 $2n - 2$ 的精确性，证明
        \[
        w_i \;=\;
        \Bigl(\,\sum_{k=0}^{n-1}
        \frac{p_k(t_i)^2}{h_k}\Bigr)^{\!-1} :
        \]
        每个权都是其节点处的
        \emph{克里斯托费尔函数}之值 --- 权的正性
        （第 10 题）再次体现，此时已有精确公式。
        验证其恢复 $w_1 = w_2 = 1$，此时 $n =
        2$、$I = \intcc{-1}1$、$w = 1$。
  \item （一个积分验证一切）对切比雪夫权
        及 $n = 3$ 的节点，计算等式两端：
        \[
        \int_{-1}^{1}\frac{t^6}{\sqrt{1 - t^2}}\,\dd t
        = \frac{5\pi}{16},
        \qquad
        Q_3(t^6) = \frac{9\pi}{32},
        \]
        由此求积误差恰为 $\frac{\pi}{32}$；
        再验证第 23 题的误差公式
        恰好预测出这一数值（此处 $f^{(6)} = 6!$ 为
        常数，且 $h_3 = \norm{2^{-2}T_3}_w^2 =
        \frac\pi{32}$）：理论与计算在
        最后一位数字上完全吻合。
\end{enumerate}
\end{problem}
